% Task C15 solutions
% Based on chapters/ch15.tex at main b4628dd90903b38fb9c71bf61921e79513bb5295
% SOLVED-PROBLEM-COUNT: 38

\chapter*{第15章\quad Fourier 级数习题解答}
\addcontentsline{toc}{chapter}{第15章\quad Fourier 级数习题解答}
\section*{§15.1 Fourier 系数：练习题}

\paragraph{练习题 1.}
\begin{xsolution}
(1) 利用恒等式
\[
  \sin^3x=\frac{3\sin x-\sin3x}{4},\qquad
  \cos^4x=\frac{3+4\cos2x+\cos4x}{8},
\]
得到
\[
  \sin^3x+\cos^4x
  =\frac38+\frac12\cos2x+\frac18\cos4x
   +\frac34\sin x-\frac14\sin3x.
\]
右边本身就是有限三角多项式，故即为所求 Fourier 级数。

(2) 在 $(-\pi,\pi)$ 上已有
\[
  x\sim2\sum_{n=1}^{\infty}\frac{(-1)^{n-1}}n\sin nx,
\]
\[
  x^2\sim\frac{\pi^2}{3}+4\sum_{n=1}^{\infty}\frac{(-1)^n}{n^2}\cos nx,
\]
以及
\[
  x^3\sim\sum_{n=1}^{\infty}(-1)^{n-1}
  \left(\frac{2\pi^2}{n}-\frac{12}{n^3}\right)\sin nx.
\]
因此线性组合给出
\begin{align*}
  ax^3+bx^2+cx+d
  &\sim \frac{b\pi^2}{3}+d
   +4b\sum_{n=1}^{\infty}\frac{(-1)^n}{n^2}\cos nx\\
  &\quad+
  \sum_{n=1}^{\infty}(-1)^{n-1}
  \left(\frac{2a\pi^2+2c}{n}-\frac{12a}{n^3}\right)\sin nx.
\end{align*}
该函数在 $(-\pi,\pi)$ 内光滑，所以 Fourier 级数在每个 $x\in(-\pi,\pi)$ 收敛于原函数；在端点处收敛于周期延拓的左右极限平均值。
\end{xsolution}

\paragraph{练习题 2.}
\begin{xsolution}
要使展开式只含奇数阶正弦项，先作奇延拓以消去全部余弦项，再令正弦系数中的偶数项消失。

在 $(0,\pi)$ 上定义
\[
  F(x)=
  \begin{cases}
    f(x),&0<x<\pi/2,\\
    f(\pi-x),&\pi/2<x<\pi,
  \end{cases}
\]
并取端点值任意，例如 $F(0)=F(\pi/2)=F(\pi)=0$；再在 $(-\pi,0)$ 上令
\[
  F(-x)=-F(x).
\]
于是 $F$ 为奇函数，所以 $a_n=0$。又对偶数指标 $2m$，
\begin{align*}
  b_{2m}
  &=\frac2\pi\int_0^\pi F(x)\sin2mx\,\dd x\\
  &=\frac2\pi\int_0^{\pi/2}
  \bigl[F(x)\sin2mx+F(\pi-x)\sin(2m(\pi-x))\bigr]\,\dd x=0,
\end{align*}
因为 $F(\pi-x)=F(x)$ 而
$\sin(2m(\pi-x))=-\sin2mx$。故 Fourier 级数只有
\[
  \sum_{m=1}^{\infty}b_{2m-1}\sin(2m-1)x
\]
中的项。
\end{xsolution}

\paragraph{练习题 3.}
% CHECK-SOLUTION: 原题把标准部分和 S_n=a_0/2+\sum_{k=1}^n(...) 与含 \sin(2nt) 的积分公式并列，索引不一致；扫描页 lower/page-0098.jpg 已核对。以下给出按标准部分和定义可直接验证的正确关系。
\begin{xsolution}
函数为奇函数，故 $a_n=0$。对 $k\ge1$，
\[
  b_k=\frac2\pi\int_0^\pi c\sin kx\,\dd x
  =\frac{2c}{\pi k}\bigl(1-(-1)^k\bigr).
\]
于是
\[
  b_{2j}=0,\qquad b_{2j-1}=\frac{4c}{\pi(2j-1)}.
\]
因此标准 Fourier 部分和满足
\[
  S_{2m-1}(x)=S_{2m}(x)
  =\frac{4c}{\pi}\sum_{j=1}^{m}\frac{\sin(2j-1)x}{2j-1}.
\]
另一方面有恒等式
\[
  \frac{\sin2mt}{\sin t}=2\sum_{j=1}^{m}\cos(2j-1)t.
\]
从 $0$ 到 $x$ 积分即得
\[
  \frac{2c}{\pi}\int_0^x\frac{\sin2mt}{\sin t}\,\dd t
  =\frac{4c}{\pi}\sum_{j=1}^{m}\frac{\sin(2j-1)x}{2j-1}.
\]
故可核验的关系为
\[
  \boxed{\displaystyle
  S_{2m-1}(x)=S_{2m}(x)
  =\frac{2c}{\pi}\int_0^x\frac{\sin2mt}{\sin t}\,\dd t }.
\]
原题若把右端的 $m$ 仍记作 $n$，则左端应相应写成 $S_{2n-1}$ 或 $S_{2n}$。
\end{xsolution}

\paragraph{练习题 4.}
\begin{xsolution}
由分部积分，
\begin{align*}
  a_n
  &=\frac1\pi\int_{-\pi}^{\pi}f(x)\cos nx\,\dd x
   =-\frac1{n\pi}\int_{-\pi}^{\pi}f'(x)\sin nx\,\dd x.
\end{align*}
因为 $f'$ 连续，Riemann 引理给出
\[
  \int_{-\pi}^{\pi}f'(x)\sin nx\,\dd x=\littleo(1),
\]
故
\[
  a_n=\littleo\!\left(\frac1n\right).
\]
同理
\begin{align*}
  b_n
  &=\frac1\pi\int_{-\pi}^{\pi}f(x)\sin nx\,\dd x\\
  &=-\frac{(-1)^n}{n\pi}\bigl[f(\pi)-f(-\pi)\bigr]
    +\frac1{n\pi}\int_{-\pi}^{\pi}f'(x)\cos nx\,\dd x.
\end{align*}
第二项为 $\littleo(1/n)$，第一项为 $\bigO(1/n)$，所以
\[
  b_n=\bigO\!\left(\frac1n\right).
\]
若再有 $f(\pi)=f(-\pi)$，边界项消失，从而
\[
  b_n=\littleo\!\left(\frac1n\right).
\]
\end{xsolution}

\paragraph{练习题 5.}
\begin{xsolution}
设 $(-\pi,\pi)$ 被有限个点分成若干单调区间 $I_j=[\alpha_j,\beta_j]$，并设
$\abs{f(x)}\le M$。只需在每个 $I_j$ 上估计振荡积分。

若 $f$ 在 $I_j$ 上单调，由积分第二中值定理，存在 $\xi_j\in I_j$，使
\[
  \int_{\alpha_j}^{\beta_j}f(x)\cos nx\,\dd x
  =f(\alpha_j+)\int_{\alpha_j}^{\xi_j}\cos nx\,\dd x
   +f(\beta_j-)\int_{\xi_j}^{\beta_j}\cos nx\,\dd x.
\]
每个余弦积分的绝对值不超过 $2/n$，故
\[
  \left|\int_{\alpha_j}^{\beta_j}f(x)\cos nx\,\dd x\right|
  \le\frac{4M}{n}.
\]
单调区间数有限，求和即得
\[
  a_n=\bigO\!\left(\frac1n\right).
\]
对正弦积分，在每个 $I_j$ 上再次用积分第二中值定理，可写成两个形如
\[
  f(\alpha_j+)\int_{\alpha_j}^{\eta_j}\sin nx\,\dd x,
  \qquad
  f(\beta_j-)\int_{\eta_j}^{\beta_j}\sin nx\,\dd x
\]
的项；每个正弦积分的绝对值也不超过 $2/n$。因此逐段求和后
\[
  b_n=\bigO\!\left(\frac1n\right).
\]
\end{xsolution}

\paragraph{练习题 6.}
\begin{xsolution}
设
\[
  f(x)\sim\frac{a_0}{2}+\sum_{n=1}^{\infty}(a_n\cos nx+b_n\sin nx),
\]
并约定 $b_0=0$。令 $g(x)=f(x)\sin x$。利用积化和差公式，对 $m\ge1$ 有
\begin{align*}
  A_m
  &=\frac1\pi\int_{-\pi}^{\pi}g(x)\cos mx\,\dd x\\
  &=\frac1{2\pi}\int_{-\pi}^{\pi}f(x)
    \bigl[\sin(m+1)x-\sin(m-1)x\bigr]\,\dd x\\
  &=\frac{b_{m+1}-b_{m-1}}2,
\end{align*}
而
\begin{align*}
  B_m
  &=\frac1\pi\int_{-\pi}^{\pi}g(x)\sin mx\,\dd x\\
  &=\frac1{2\pi}\int_{-\pi}^{\pi}f(x)
    \bigl[\cos(m-1)x-\cos(m+1)x\bigr]\,\dd x\\
  &=\frac{a_{m-1}-a_{m+1}}2,
\end{align*}
其中 $m=1$ 时 $a_{m-1}=a_0$。此外
\[
  A_0=\frac1\pi\int_{-\pi}^{\pi}f(x)\sin x\,\dd x=b_1.
\]
另一方面，形式上逐项乘以 $\sin x$ 后
\begin{align*}
  \sin x\left[\frac{a_0}{2}+\sum_{n\ge1}(a_n\cos nx+b_n\sin nx)\right]
\end{align*}
用同样的积化和差公式整理，恰好得到常数项 $A_0/2$、余弦系数 $A_m$ 与正弦系数 $B_m$。因此所得三角级数正是 $f(x)\sin x$ 的 Fourier 级数。
\end{xsolution}

\paragraph{练习题 7.}
\begin{xsolution}
令
\[
  E_n(a_1,\ldots,a_n)
  =\int_a^b\left[f(x)-\sum_{k=1}^{n}a_ke_k(x)\right]^2\,\dd x.
\]
利用规范正交性，展开得
\begin{align*}
  E_n
  &=\int_a^bf^2(x)\,\dd x
    -2\sum_{k=1}^{n}a_kc_k+\sum_{k=1}^{n}a_k^2\\
  &=\int_a^bf^2(x)\,\dd x-\sum_{k=1}^{n}c_k^2
    +\sum_{k=1}^{n}(a_k-c_k)^2.
\end{align*}
因为 $E_n\ge0$，取 $a_k=c_k$ 即得 Bessel 不等式
\[
  \sum_{k=1}^{n}c_k^2\le\int_a^bf^2(x)\,\dd x.
\]
左端随 $n$ 单调增加且有统一上界，故
\[
  \sum_{n=1}^{\infty}c_n^2
\]
收敛。

同时上面的配方公式说明，当 $n$ 固定时，$E_n$ 的唯一最小值在
\[
  \boxed{a_k=c_k\quad(k=1,\ldots,n)}
\]
处取得。
\end{xsolution}

\paragraph{练习题 8.}
\begin{xsolution}
令
\[
  G(x)=\int_0^xg(t)\,\dd t.
\]
由于 $g$ 的周期为 $1$ 且 $\int_0^1g=0$，有 $G(x+1)=G(x)$，所以 $G$ 连续且有界。又
\[
  \frac{\dd}{\dd x}G(nx)=n g(nx).
\]
因此分部积分给出
\begin{align*}
  a_n
  &=\frac1n\int_0^1f(x)\,\dd[G(nx)]\\
  &=\frac1n\bigl[f(x)G(nx)\bigr]_0^1
    -\frac1n\int_0^1f'(x)G(nx)\,\dd x.
\end{align*}
取 $G(0)=0$，则 $G(n)=G(0)=0$，边界项为零。若
$\norm{G}_\infty=M$，则
\[
  \abs{a_n}\le\frac{M}{n}\int_0^1\abs{f'(x)}\,\dd x=\bigO\!\left(\frac1n\right).
\]
于是
\[
  \sum_{n=1}^{\infty}a_n^2
\]
由与 $\sum1/n^2$ 比较而收敛。
\end{xsolution}

\section*{§15.2 Fourier 级数的收敛性：练习题}

\paragraph{练习题 1.}
\begin{xsolution}
作偶延拓，故只需求余弦系数。由直接积分
\[
  a_0=\frac2\pi\int_0^\pi f(x)\,\dd x=\frac\pi2,
\]
且对 $n\ge1$，
\begin{align*}
  a_n
  &=\frac2\pi\left[
    \int_0^{\pi/2}x\cos nx\,\dd x
    +\int_{\pi/2}^{\pi}\left(x-\frac\pi2\right)\cos nx\,\dd x
    \right]\\
  &=\frac{\sin(n\pi/2)}n
    +\frac{2[(-1)^n-1]}{\pi n^2}.
\end{align*}
因此偶数项全为零，而
\[
  a_{2m-1}
  =\frac{(-1)^{m-1}}{2m-1}
   -\frac4{\pi(2m-1)^2}.
\]
故余弦展开式为
\[
  f(x)\sim\frac\pi4+
  \sum_{m=1}^{\infty}
  \left[
    \frac{(-1)^{m-1}}{2m-1}
    -\frac4{\pi(2m-1)^2}
  \right]\cos(2m-1)x.
\]
在 $x=0$ 处周期偶延拓连续，故级数和为 $f(0)=0$。于是
\[
  0=\frac\pi4+\sum_{m=1}^{\infty}\frac{(-1)^{m-1}}{2m-1}
  -\frac4\pi\sum_{m=1}^{\infty}\frac1{(2m-1)^2}.
\]
由已知 $\sum_{n\ge1}1/n^2=\pi^2/6$，
\[
  \sum_{m=1}^{\infty}\frac1{(2m-1)^2}
  =\frac{\pi^2}{6}-\frac14\frac{\pi^2}{6}=\frac{\pi^2}{8}.
\]
代回即得
\[
  \boxed{\displaystyle
  \sum_{m=1}^{\infty}\frac{(-1)^{m-1}}{2m-1}=\frac\pi4 }.
\]
\end{xsolution}

\paragraph{练习题 2.}
\begin{xsolution}
在 $[0,a]$ 上作半区间展开，基本频率为 $\pi/a$。

(1) 余弦级数的系数为
\[
  a_0=\frac2a\int_0^af(x)\,\dd x=\frac a2,
\]
以及
\[
  a_n=\frac2a\int_0^af(x)\cos\frac{n\pi x}{a}\,\dd x
  =\frac{2a}{\pi^2n^2}
  \left[(-1)^{n+1}+2\cos\frac{n\pi}{2}-1\right].
\]
只有 $n=4k+2$ 时非零，此时
\[
  a_{4k+2}=-\frac{2a}{\pi^2(2k+1)^2}.
\]
故
\[
  \boxed{\displaystyle
  f(x)=\frac{a}{4}-
  \frac{2a}{\pi^2}
  \sum_{k=0}^{\infty}
  \frac{\cos\frac{2(2k+1)\pi x}{a}}{(2k+1)^2}}
  \qquad(0\le x\le a).
\]

(2) 正弦级数的系数为
\[
  b_n=\frac2a\int_0^af(x)\sin\frac{n\pi x}{a}\,\dd x
  =\frac{4a}{\pi^2n^2}\sin\frac{n\pi}{2}.
\]
故只有奇数项，
\[
  \boxed{\displaystyle
  f(x)=\frac{4a}{\pi^2}
  \sum_{k=0}^{\infty}
  \frac{(-1)^k}{(2k+1)^2}
  \sin\frac{(2k+1)\pi x}{a}}
  \qquad(0\le x\le a).
\]
函数连续且分段光滑，两式在开区间内均收敛于 $f$；这里 $f(0)=f(a)=0$，正弦级数在端点也给出正确值。
\end{xsolution}

\paragraph{练习题 3.}
\begin{xsolution}
设 $\alpha\notin\ZZ$。函数 $\cos\alpha x$ 在 $[-\pi,\pi]$ 上为偶函数，其 Fourier 系数为
\[
  a_0=\frac{2\sin\pi\alpha}{\pi\alpha},
\]
以及
\begin{align*}
  a_n
  &=\frac2\pi\int_0^\pi\cos\alpha x\cos nx\,\dd x\\
  &=\frac{2\alpha(-1)^n\sin\pi\alpha}{\pi(\alpha^2-n^2)},\qquad n\ge1.
\end{align*}
因此
\[
  \cos\alpha x
  =\frac{\sin\pi\alpha}{\pi\alpha}
  +\frac{2\alpha\sin\pi\alpha}{\pi}
  \sum_{n=1}^{\infty}\frac{(-1)^n\cos nx}{\alpha^2-n^2}.
\]
令 $x=0$ 并除以 $\sin\pi\alpha$，得到
\[
  \csc(\pi\alpha)
  =\frac1{\pi\alpha}
  +\frac{2\alpha}{\pi}
  \sum_{n=1}^{\infty}\frac{(-1)^n}{\alpha^2-n^2}.
\]
令 $x=\pi$，得到
\[
  \cot(\pi\alpha)
  =\frac1{\pi\alpha}
  +\frac{2\alpha}{\pi}
  \sum_{n=1}^{\infty}\frac1{\alpha^2-n^2}.
\]
把 $z=\pi\alpha$ 代入即得
\[
  \boxed{\displaystyle
  \cot z=\frac1z+\sum_{n=1}^{\infty}
  \left(\frac1{z-n\pi}+\frac1{z+n\pi}\right)
  =\frac1z+\sum_{n=1}^{\infty}\frac{2z}{z^2-n^2\pi^2}},
\]
\[
  \boxed{\displaystyle
  \csc z=\frac1z+\sum_{n=1}^{\infty}(-1)^n
  \left(\frac1{z-n\pi}+\frac1{z+n\pi}\right)
  =\frac1z+\sum_{n=1}^{\infty}(-1)^n\frac{2z}{z^2-n^2\pi^2}},
\]
其中 $z\notin\pi\ZZ$。

最后取 $z=\pi a$，由余切展开式
\[
  \cot\pi a=\frac1{\pi a}+\frac{2a}{\pi}
  \sum_{n=1}^{\infty}\frac1{a^2-n^2},
\]
所以对 $a\notin\ZZ$ 且 $a\ne0$，
\[
  \boxed{\displaystyle
  \sum_{n=1}^{\infty}\frac1{n^2-a^2}
  =\frac1{2a^2}-\frac{\pi}{2a}\cot\pi a }.
\]
\end{xsolution}

\paragraph{练习题 4.}
\begin{xsolution}
(1) $f(x)=\abs{\sin x}$ 为偶函数。于是 $b_n=0$，
\[
  a_0=\frac2\pi\int_0^\pi\sin x\,\dd x=\frac4\pi.
\]
对 $n\ge1$，
\[
  a_n=\frac2\pi\int_0^\pi\sin x\cos nx\,\dd x.
\]
当 $n$ 为奇数时 $a_n=0$；令 $n=2m$ 时，
\[
  a_{2m}=-\frac4{\pi(4m^2-1)}.
\]
故
\[
  \boxed{\displaystyle
  \abs{\sin x}=\frac{2}{\pi}-
  \frac{4}{\pi}
  \sum_{m=1}^{\infty}\frac{\cos2mx}{4m^2-1}}.
\]
该级数绝对一致收敛。

(2) 记
\[
  f(x)=\begin{cases}ax,&-\pi\le x<0,\\ bx,&0\le x\le\pi.\end{cases}
\]
直接计算得
\[
  a_0=\frac1\pi\int_{-\pi}^{\pi}f(x)\,\dd x
  =\frac\pi2(b-a),
\]
\[
  a_n=\frac{b-a}{\pi n^2}\bigl[(-1)^n-1\bigr],
\qquad
  b_n=\frac{(a+b)(-1)^{n+1}}n.
\]
因此
\[
  \boxed{\displaystyle
  f(x)\sim\frac{\pi(b-a)}4
  +\sum_{n=1}^{\infty}
   \frac{(b-a)[(-1)^n-1]}{\pi n^2}\cos nx
  +(a+b)\sum_{n=1}^{\infty}\frac{(-1)^{n+1}}n\sin nx }.
\]
在 $(-\pi,\pi)$ 内级数收敛于 $f(x)$；在端点 $\pm\pi$ 处收敛于周期延拓左右极限的平均值
$\pi(b-a)/2$。
\end{xsolution}

\paragraph{练习题 5.}
\begin{xsolution}
(1) 在 $[0,\pi]$ 作正弦展开。对 $n\ge1$，
\begin{align*}
  b_n
  &=\frac2\pi\int_0^\pi\ee^{-2x}\sin nx\,\dd x\\
  &=\frac{2n}{\pi(n^2+4)}
    \left[1-(-1)^n\ee^{-2\pi}\right].
\end{align*}
故
\[
  \boxed{\displaystyle
  \ee^{-2x}\sim
  \sum_{n=1}^{\infty}
  \frac{2n[1-(-1)^n\ee^{-2\pi}]}{\pi(n^2+4)}\sin nx,
  \qquad0<x<\pi }.
\]
由于作奇延拓后在 $x=0$ 与 $x=\pi$ 都有第一类间断点，正弦级数在两端点的和均为左右极限的平均值 $0$；在 $0<x<\pi$ 收敛于原函数。

(2) 区间长度为 $2$，故正弦基为 $\sin(n\pi x/2)$。有
\[
  b_n=\int_0^2f(x)\sin\frac{n\pi x}{2}\,\dd x
  =\int_0^1\cos\frac{\pi x}{2}\sin\frac{n\pi x}{2}\,\dd x.
\]
当 $n=1$ 时
\[
  b_1=\frac1\pi.
\]
当 $n>1$ 时由积化和差公式
\[
  b_n=\frac{2[n-\sin(n\pi/2)]}{\pi(n^2-1)}.
\]
因此
\[
  \boxed{\displaystyle
  f(x)\sim\frac1\pi\sin\frac{\pi x}{2}
  +\sum_{n=2}^{\infty}
  \frac{2[n-\sin(n\pi/2)]}{\pi(n^2-1)}
  \sin\frac{n\pi x}{2} }.
\]
在 $0<x<2$ 的连续点处收敛于 $f$。在端点处，正弦级数的和均为 $0$；因此在 $x=2$ 与题设函数值一致，而在 $x=0$ 收敛于奇延拓左右极限的平均值 $0$，不等于题设的 $f(0)=1$。
\end{xsolution}

\paragraph{练习题 6.}
\begin{xsolution}
(1) 在 $[0,\pi]$ 作余弦展开。直接积分得
\[
  a_0=\frac2\pi\int_0^\pi x(\pi-x)\,\dd x=\frac{\pi^2}{3},
\]
以及
\[
  a_n=\frac2\pi\int_0^\pi x(\pi-x)\cos nx\,\dd x
  =\frac{2[(-1)^{n+1}-1]}{n^2}.
\]
所以只有偶数项，且 $a_{2m}=-1/m^2$。于是
\[
  \boxed{\displaystyle
  x(\pi-x)=\frac{\pi^2}{6}-
  \sum_{m=1}^{\infty}\frac{\cos2mx}{m^2},
  \qquad0\le x\le\pi }.
\]

(2) 这里区间为 $[0,\pi/2]$，余弦基为 $\cos2nx$。有
\[
  a_0=\frac4\pi\left(
  \int_0^{\pi/4}\sin2x\,\dd x+
  \int_{\pi/4}^{\pi/2}1\,\dd x\right)
  =\frac{\pi+2}{\pi}.
\]
对 $n=1$，
\[
  a_1=-\frac1\pi.
\]
对 $n>1$，
\[
  a_n=\frac4\pi\left(
  \int_0^{\pi/4}\sin2x\cos2nx\,\dd x+
  \int_{\pi/4}^{\pi/2}\cos2nx\,\dd x\right)
  =\frac{2[-n+\sin(n\pi/2)]}{\pi n(n^2-1)}.
\]
故
\[
  \boxed{\displaystyle
  f(x)=\frac{\pi+2}{2\pi}-\frac1\pi\cos2x
  +\sum_{n=2}^{\infty}
  \frac{2[-n+\sin(n\pi/2)]}{\pi n(n^2-1)}\cos2nx }.
\]
原函数连续且分段光滑，故在 $[0,\pi/2]$ 上按偶延拓的 Fourier 余弦级数收敛于相应连续值。
\end{xsolution}

\paragraph{练习题 7.}
\begin{xsolution}
函数为奇函数，故只有正弦项。对 $n\ge1$，
\begin{align*}
  b_n
  &=\frac2\pi\int_0^\pi(\pi-x)\sin nx\,\dd x
  =\frac2n.
\end{align*}
所以
\[
  \boxed{\displaystyle
  f(x)\sim2\sum_{n=1}^{\infty}\frac{\sin nx}{n}}.
\]
函数的周期延拓在 $x=0$ 有第一类间断点，左右极限分别为 $-\pi$ 与 $\pi$，而 $f(0)=0$ 正好是平均值；在其余点分段光滑。因此 Fourier 级数在 $(-\pi,\pi]$ 每一点都收敛于题设的 $f$。

但该收敛不可能一致，因为每个部分和都是连续函数，而 $f$ 在 $x=0$ 不连续；连续函数列的一致极限必连续。故级数不一致收敛。
\end{xsolution}

\paragraph{练习题 8.}
\begin{xsolution}
先取连续的 $2\pi$ 周期函数 $g$，使
\[
  \int_{-\pi}^{\pi}\abs{f(x)-g(x)}\,\dd x<\frac\varepsilon2.
\]
这样的 $g$ 总能取得：先用连续函数在 $L^1$ 意义下逼近 $f$，再只在端点的任意小邻域内作线性连接，使两端值相等；由于修改区间可以任意短，增加的 $L^1$ 误差也可任意小。记 $g$ 的 Fourier 部分和为 $S_0,S_1,\ldots$，并令其 Fejér 平均为
\[
  \sigma_n(x)=\frac{S_0(x)+S_1(x)+\cdots+S_n(x)}{n+1}.
\]
每个 $\sigma_n$ 都是三角多项式。由一致收敛形式的 Fejér 定理，
\[
  \sigma_n\to g
\]
在 $[-\pi,\pi]$ 上一致收敛，故可取 $n$ 充分大，使
\[
  \max_{[-\pi,\pi]}\abs{g(x)-\sigma_n(x)}
  <\frac{\varepsilon}{4\pi}.
\]
于是
\begin{align*}
  \int_{-\pi}^{\pi}\abs{f(x)-\sigma_n(x)}\,\dd x
  &\le\int_{-\pi}^{\pi}\abs{f-g}\,\dd x
   +\int_{-\pi}^{\pi}\abs{g-\sigma_n}\,\dd x\\
  &<\frac\varepsilon2+2\pi\cdot\frac{\varepsilon}{4\pi}
  =\varepsilon.
\end{align*}
取 $P_n=\sigma_n$ 即得所求三角多项式。
\end{xsolution}

\paragraph{练习题 9.}
\begin{xsolution}
设右边三角级数一致收敛于 $g$。由于每个部分和连续，$g$ 连续。又由命题 15.1.1，一致收敛的三角级数必是其和函数 $g$ 的 Fourier 级数，因此 $g$ 与 $f$ 有完全相同的 Fourier 系数。

函数 $f-g$ 连续且全部 Fourier 系数均为 $0$。由 Fourier 级数惟一性定理（命题 15.2.3），
\[
  f-g\equiv0.
\]
故一致收敛时其和函数必为 $f$。
\end{xsolution}

\paragraph{练习题 10.}
\begin{xsolution}
函数为偶函数，故 $b_n=0$，并有
\[
  a_0=\frac2\pi\int_0^af(x)\,\dd x=\frac{2a}{\pi},
\qquad
  a_n=\frac2\pi\int_0^a\cos nx\,\dd x
  =\frac{2\sin na}{\pi n}.
\]
Parseval 等式给出
\[
  \frac12\left(\frac{2a}{\pi}\right)^2
  +\sum_{n=1}^{\infty}\left(\frac{2\sin na}{\pi n}\right)^2
  =\frac1\pi\int_{-\pi}^{\pi}f^2(x)\,\dd x
  =\frac{2a}{\pi}.
\]
整理得
\[
  \boxed{\displaystyle
  \sum_{n=1}^{\infty}\frac{\sin^2na}{n^2}
  =\frac{a(\pi-a)}2}.
\]
再由
\[
  \sum_{n=1}^{\infty}\frac1{n^2}=\frac{\pi^2}{6}
\]
和 $\cos^2na=1-\sin^2na$，得到
\[
  \boxed{\displaystyle
  \sum_{n=1}^{\infty}\frac{\cos^2na}{n^2}
  =\frac{\pi^2}{6}-\frac{a(\pi-a)}2}.
\]
\end{xsolution}

\section*{§15.3 参考题}

\paragraph{参考题 1.}
\begin{xsolution}
(1) 把 $[0,2\pi]$ 分成 $n$ 个长度为 $2\pi/n$ 的周期小区间。对 $j=0,1,\ldots,n-1$，令
\[
  I_j=\int_{2j\pi/n}^{2(j+1)\pi/n}f(x)\sin nx\,\dd x.
\]
把后一半区间平移 $\pi/n$，得到
\[
  I_j=\int_0^{\pi/n}
  \left[f\left(\frac{2j\pi}{n}+t\right)
  -f\left(\frac{(2j+1)\pi}{n}+t\right)\right]\sin nt\,\dd t.
\]
由于 $f$ 单调减少，方括号内非负；而 $0\le t\le\pi/n$ 时 $\sin nt\ge0$。故 $I_j\ge0$。求和即得
\[
  \boxed{\displaystyle
  \int_0^{2\pi}f(x)\sin nx\,\dd x\ge0}.
\]
若端点为广义积分，只需先在内闭区间作上述运算，再令端点趋近即可。

(2) 由 $f'$ 可积和绝对可积，分部积分合法，而且 $\sin0=\sin2n\pi=0$，故
\[
  \int_0^{2\pi}f(x)\cos nx\,\dd x
  =-\frac1n\int_0^{2\pi}f'(x)\sin nx\,\dd x.
\]
因为 $f'$ 单调增加，所以 $-f'$ 单调减少。由 (1)，
\[
  \int_0^{2\pi}[-f'(x)]\sin nx\,\dd x\ge0.
\]
于是
\[
  \boxed{\displaystyle
  \int_0^{2\pi}f(x)\cos nx\,\dd x\ge0}.
\]
\end{xsolution}

\paragraph{参考题 2.}
\begin{xsolution}
端点的单点取值不影响积分，因此必要时可把 $f(0),f(2\pi)$ 改为相应的单侧极限。这样可把 $f$ 看成闭区间上的连续下凸函数。下凸函数在任意内闭区间上绝对连续，其导函数几乎处处存在且有一个单调增加的代表；令端点趋近即可得到分部积分公式
\[
  \int_0^{2\pi}f(x)\cos nx\,\dd x
  =-\frac1n\int_0^{2\pi}f'(x)\sin nx\,\dd x.
\]
这里 $-f'$ 单调减少。参考题 1(1) 的证明只用到了单调性，故同样适用于 $-f'$，于是
\[
  \int_0^{2\pi}[-f'(x)]\sin nx\,\dd x\ge0.
\]
因此
\[
  \boxed{\displaystyle
  \int_0^{2\pi}f(x)\cos nx\,\dd x\ge0,\qquad n\ge1 }.
\]
\end{xsolution}

\paragraph{参考题 3.}
\begin{xsolution}
(1) 由周期性作变量平移可知
\[
  F_h(x+2\pi)=F_h(x).
\]
又因 $f$ 连续，变上、下限积分求导给出
\[
  F_h'(x)=\frac{f(x+h)-f(x-h)}{2h},
\]
右边连续，所以 $F_h$ 是周期 $2\pi$ 的连续可微函数。

(2) 周期连续函数 $f$ 在 $\RR$ 上一致连续。给定 $\varepsilon>0$，取 $\delta>0$ 使
$\abs{u}<\delta$ 时
\[
  \abs{f(x+u)-f(x)}<\varepsilon
\]
对所有 $x$ 一致成立。若 $0<h<\delta$，则
\begin{align*}
  \abs{F_h(x)-f(x)}
  &=\left|\frac1{2h}\int_{-h}^{h}[f(x+u)-f(x)]\,\dd u\right|\\
  &\le\frac1{2h}\int_{-h}^{h}\abs{f(x+u)-f(x)}\,\dd u<\varepsilon.
\end{align*}
故 $F_h\to f$ 一致。

(3) 给定 $\varepsilon>0$，先由 (2) 取 $h$ 使
\[
  \norm{F_h-f}_\infty<\frac\varepsilon2.
\]
由 (1)，$F_h\in C^1$，其导函数连续，故满足命题 15.2.8 的条件；于是 $F_h$ 的 Fourier 级数绝对一致收敛于 $F_h$。取其中一个足够高阶的部分和三角多项式 $P_N$，使
\[
  \norm{F_h-P_N}_\infty<\frac\varepsilon2.
\]
于是
\[
  \norm{f-P_N}_\infty<\varepsilon,
\]
这就重新证明了周期函数形式的 Weierstrass 第二逼近定理。

(4) 设 $f$ 的 Fourier 系数为 $a_0,a_n,b_n$。对 $n\ge1$，由 Fubini 换序与周期平移，
\begin{align*}
  A_n
  &=\frac1\pi\int_{-\pi}^{\pi}F_h(x)\cos nx\,\dd x\\
  &=\frac1{2h}\int_{-h}^{h}
  \frac1\pi\int_{-\pi}^{\pi}f(x+u)\cos nx\,\dd x\,\dd u\\
  &=\frac1{2h}\int_{-h}^{h}
  (a_n\cos nu+b_n\sin nu)\,\dd u\\
  &=a_n\frac{\sin nh}{nh}.
\end{align*}
对于正弦系数，同样作 $y=x+u$，内层积分为
\[
  \int_{-\pi}^{\pi}f(x+u)\sin nx\,\dd x
  =\pi(b_n\cos nu-a_n\sin nu).
\]
于是对称区间 $[-h,h]$ 上的奇函数项积分为零，从而
\[
  B_n=b_n\frac{\sin nh}{nh}.
\]
常数项则直接由周期平移得到 $A_0=a_0$。
因此
\[
  \boxed{\displaystyle
  F_h(x)\sim\frac{a_0}{2}
  +\sum_{n=1}^{\infty}\frac{\sin nh}{nh}
  (a_n\cos nx+b_n\sin nx)}.
\]
\end{xsolution}

\paragraph{参考题 4.}
\begin{xsolution}
(1) 先求 $F$ 的 Fourier 系数。对 $m\ge1$，
\begin{align*}
  A_m
  &=\frac1\pi\int_{-\pi}^{\pi}F(x)\cos mx\,\dd x\\
  &=\frac1{\pi^2}\int_{-\pi}^{\pi}f(t)
    \left[\int_{-\pi}^{\pi}f(x+t)\cos mx\,\dd x\right]\dd t.
\end{align*}
令 $y=x+t$，利用周期性，内层积分等于
\[
  \pi(a_m\cos mt+b_m\sin mt).
\]
因此
\[
  A_m=a_m^2+b_m^2.
\]
对正弦系数，变量代换后内层积分为
\[
  \int_{-\pi}^{\pi}f(x+t)\sin mx\,\dd x
  =\pi(b_m\cos mt-a_m\sin mt).
\]
因此
\begin{align*}
  B_m
  &=\frac1\pi\int_{-\pi}^{\pi}
    f(t)(b_m\cos mt-a_m\sin mt)\,\dd t\\
  &=b_ma_m-a_mb_m=0.
\end{align*}
常数项为
\[
  A_0=a_0^2.
\]
故
\[
  \boxed{\displaystyle
  F(x)\sim\frac{a_0^2}{2}
  +\sum_{n=1}^{\infty}(a_n^2+b_n^2)\cos nx }.
\]

(2) 连续函数 $f$ 平方可积，由 Bessel 不等式
\[
  \sum_{n=1}^{\infty}(a_n^2+b_n^2)<\infty.
\]
所以由 Weierstrass M-判别法，$F$ 的上述 Fourier 级数绝对一致收敛。

(3) 由定义，$F$ 连续；由 (2) 的一致收敛和命题 15.1.1，右边级数的和就是 $F$。令 $x=0$，得到
\[
  \frac1\pi\int_{-\pi}^{\pi}f^2(t)\,\dd t
  =F(0)=\frac{a_0^2}{2}+\sum_{n=1}^{\infty}(a_n^2+b_n^2),
\]
即 Parseval 等式。
\end{xsolution}

\paragraph{参考题 5.}
\begin{xsolution}
(1) 记
\[
  S_n(x)=\sum_{k=1}^{n}\frac{\sin kx}{k}.
\]
由题给恒等式和 Dirichlet 核公式，
\begin{align*}
  S_n(x)
  &=\int_0^x\sum_{k=1}^{n}\cos kt\,\dd t\\
  &=\frac12\int_0^x
  \frac{\sin(n+\frac12)t}{\sin(t/2)}\,\dd t-\frac x2.
\end{align*}
当 $0<x<2\pi$ 时，函数
\[
  q(t)=\frac{t}{2\sin(t/2)}
\]
在 $[0,x]$ 上连续且 $q(0)=1$。故由 Dirichlet 积分极限
\[
  \lim_{\lambda\to\infty}\int_0^xq(t)\frac{\sin\lambda t}{t}\,\dd t=\frac\pi2,
\]
可得
\[
  \boxed{\displaystyle
  \sum_{k=1}^{\infty}\frac{\sin kx}{k}=\frac{\pi-x}{2},
  \qquad0<x<2\pi }.
\]
在 $x\in2\pi\ZZ$ 时各项均为 $0$，和为 $0$；其余值由 $2\pi$ 周期性确定。

(2) 记上式和函数为 $S(x)$。则
\[
  \sum_{k=1}^{\infty}\frac{\sin(2k-1)x}{2k-1}
  =S(x)-\frac12S(2x).
\]
当 $0<x<\pi$ 时，
\[
  S(x)=\frac{\pi-x}{2},\qquad
  S(2x)=\frac{\pi-2x}{2},
\]
故
\[
  \boxed{\displaystyle
  \sum_{k=1}^{\infty}\frac{\sin(2k-1)x}{2k-1}=\frac\pi4,
  \qquad0<x<\pi }.
\]
当 $\pi<x<2\pi$ 时，$S(x)=(\pi-x)/2$，而 $2x-2\pi\in(0,2\pi)$，故由周期性
\[
  S(2x)=S(2x-2\pi)=\frac{3\pi-2x}{2}.
\]
代入 $S(x)-\frac12S(2x)$ 得 $-\pi/4$。在 $x\in\pi\ZZ$ 时每一项均为 $0$，故和为 $0$。
\end{xsolution}

\paragraph{参考题 6.}
\begin{xsolution}
以下均从参考题 5(1) 的
\[
  S(x)=\sum_{n=1}^{\infty}\frac{\sin nx}{n}=\frac{\pi-x}{2}
  \qquad(0<x<2\pi)
\]
出发。

(1) 偶数项为
\[
  \sum_{k=1}^{\infty}\frac{\sin2kx}{2k}=\frac12S(2x)
  =\boxed{\frac\pi4-\frac x2},\qquad0<x<\pi.
\]

(2) 用总和减去偶数项，
\[
  \sum_{k=1}^{\infty}\frac{\sin(2k-1)x}{2k-1}
  =S(x)-\frac12S(2x)=\boxed{\frac\pi4},
  \qquad0<x<\pi.
\]

(3) 对 $\abs{x}<\pi$，令 $y=x+\pi\in(0,2\pi)$。由于
$\sin n(x+\pi)=(-1)^n\sin nx$，
\[
  \sum_{n=1}^{\infty}\frac{(-1)^{n-1}}n\sin nx
  =-S(x+\pi)=\boxed{\frac x2}.
\]

(4) 对 (3) 从 $0$ 到 $x$ 逐项积分。积分后的级数由 $1/n^2$ 控制，故一致收敛，得到
\[
  \frac{x^2}{4}
  =\sum_{n=1}^{\infty}\frac{(-1)^{n-1}}{n^2}(1-\cos nx).
\]
先令 $x=\pi$，得到
\[
  2\sum_{m=1}^{\infty}\frac1{(2m-1)^2}=\frac{\pi^2}{4},
\]
故奇数平方倒数和为 $\pi^2/8$。若记 $Z=\sum1/n^2$，则
$Z=\pi^2/8+Z/4$，所以 $Z=\pi^2/6$，进而
\[
  \sum_{n=1}^{\infty}\frac{(-1)^{n-1}}{n^2}=\frac{\pi^2}{12}.
\]
代回前式即得
\[
  \boxed{\displaystyle
  x^2=\frac{\pi^2}{3}
  +4\sum_{n=1}^{\infty}\frac{(-1)^n}{n^2}\cos nx,
  \qquad\abs{x}<\pi }.
\]
端点按连续延拓也成立。

(5) 对 (2) 从 $0$ 到 $x$ 积分，利用
$\sum_{k\ge1}(2k-1)^{-2}=\pi^2/8$，
\[
  \frac{\pi x}{4}
  =\frac{\pi^2}{8}-
  \sum_{k=1}^{\infty}\frac{\cos(2k-1)x}{(2k-1)^2}.
\]
整理得
\[
  \boxed{\displaystyle
  x=\frac\pi2-\frac4\pi
  \sum_{k=1}^{\infty}\frac{\cos(2k-1)x}{(2k-1)^2},
  \qquad0\le x\le\pi }.
\]

(6) 对参考题 5(1) 从 $0$ 到 $x$ 积分：
\[
  \sum_{n=1}^{\infty}\frac{1-\cos nx}{n^2}
  =\frac{\pi x}{2}-\frac{x^2}{4}.
\]
再用 $\sum1/n^2=\pi^2/6$，得到
\[
  \boxed{\displaystyle
  \sum_{n=1}^{\infty}\frac{\cos nx}{n^2}
  =\frac{3x^2-6\pi x+2\pi^2}{12},
  \qquad0\le x\le2\pi }.
\]
题中所列 $0\le x\le\pi$ 是其子区间。
\end{xsolution}

\paragraph{参考题 7（Steklov（斯捷克洛夫）不等式）.}
\begin{xsolution}
先考虑条件 (1)：$\int_0^\pi f=0$。把 $f$ 作半区间余弦展开
\[
  f(x)\sim\sum_{n=1}^{\infty}a_n\cos nx,
\qquad
  a_n=\frac2\pi\int_0^\pi f(x)\cos nx\,\dd x,
\]
其中常数项因均值为零而消失。Parseval 等式给出
\[
  \int_0^\pi f^2(x)\,\dd x
  =\frac\pi2\sum_{n=1}^{\infty}a_n^2.
\]
另一方面，对 $f'$ 的正弦系数有
\[
  \frac2\pi\int_0^\pi f'(x)\sin nx\,\dd x
  =-n a_n,
\]
因为边界上的 $\sin nx$ 为零。由 Bessel 不等式，
\[
  \int_0^\pi f'^2(x)\,\dd x
  \ge\frac\pi2\sum_{n=1}^{\infty}n^2a_n^2
  \ge\frac\pi2\sum_{n=1}^{\infty}a_n^2
  =\int_0^\pi f^2(x)\,\dd x.
\]
若等号成立，则 $a_n=0$ 对一切 $n\ge2$，由余弦系完备性与连续性，
\[
  f(x)=A\cos x.
\]
反之该函数满足均值为零并取得等号。

再考虑条件 (2)：$f(0)=f(\pi)=0$。作半区间正弦展开
\[
  f(x)\sim\sum_{n=1}^{\infty}b_n\sin nx,
\qquad
  \int_0^\pi f^2(x)\,\dd x=\frac\pi2\sum_{n=1}^{\infty}b_n^2.
\]
对 $f'$ 的余弦系数，分部积分并用端点条件得到
\[
  \frac2\pi\int_0^\pi f'(x)\cos nx\,\dd x=n b_n,
\]
而常数余弦系数为 $0$。故 Bessel 不等式给出
\[
  \int_0^\pi f'^2(x)\,\dd x
  \ge\frac\pi2\sum_{n=1}^{\infty}n^2b_n^2
  \ge\int_0^\pi f^2(x)\,\dd x.
\]
等号当且仅当只有 $b_1$ 可能非零，即
\[
  f(x)=B\sin x.
\]
这正是题中两种条件各自的等号情形。
\end{xsolution}

\paragraph{参考题 8（Wirtinger 不等式）.}
\begin{xsolution}
设
\[
  f(x)\sim\frac{a_0}{2}+\sum_{n=1}^{\infty}(a_n\cos nx+b_n\sin nx).
\]
由 $\int_{-\pi}^{\pi}f=0$ 知 $a_0=0$。Parseval 等式给出
\[
  \frac1\pi\int_{-\pi}^{\pi}f^2(x)\,\dd x
  =\sum_{n=1}^{\infty}(a_n^2+b_n^2).
\]
因为 $f(-\pi)=f(\pi)$，命题 15.1.2 给出 $f'$ 的 Fourier 系数
\[
  a_n'=n b_n,\qquad b_n'=-n a_n.
\]
再对 $f'$ 用 Parseval 等式，
\[
  \frac1\pi\int_{-\pi}^{\pi}f'^2(x)\,\dd x
  =\sum_{n=1}^{\infty}n^2(a_n^2+b_n^2)
  \ge\sum_{n=1}^{\infty}(a_n^2+b_n^2).
\]
故
\[
  \boxed{\displaystyle
  \int_{-\pi}^{\pi}f'^2(x)\,\dd x
  \ge\int_{-\pi}^{\pi}f^2(x)\,\dd x }.
\]
等号成立当且仅当 $n\ge2$ 的全部系数均为 $0$，即
\[
  \boxed{f(x)=A\cos x+B\sin x}.
\]
这类函数确实满足题设条件并取得等号。
\end{xsolution}

\paragraph{参考题 9.}
\begin{xsolution}
设闭区间 $I$ 不含 $f$ 的任何间断点。因为 $f,f'$ 均分段连续，间断点只有有限个，所以可取 $\eta>0$，使每个 $x\in I$ 的 $\eta$ 邻域都不含 $f$ 的间断点。

在各间断点的小邻域内把 $f$ 用线段连接，得到一个周期 $2\pi$ 的连续分段 $C^1$ 函数 $g$，并使 $g=f$ 在 $I$ 的固定邻域上成立。于是 $g'$ 分段连续并平方可积，由命题 15.2.8，$g$ 的 Fourier 级数绝对一致收敛于 $g$。

令 $h=f-g$。则对 $x\in I$ 和 $\abs{t}<\eta$ 有
\[
  h(x+t)=h(x-t)=0.
\]
由 Dirichlet 积分表示，$h$ 的第 $n$ 个 Fourier 部分和满足
\[
  S_n^h(x)=\frac1\pi\int_{\eta}^{\pi}
  [h(x+t)+h(x-t)]
  \frac{\sin(n+\frac12)t}{2\sin(t/2)}\,\dd t.
\]
对 $x\in I$，积分中的振幅关于 $t$ 分段 $C^1$，且由于 $t\ge\eta$，分母远离 $0$；其总变差可由一个与 $x$ 无关的常数控制。对每个分段作分部积分，得到
\[
  \sup_{x\in I}\abs{S_n^h(x)}=\bigO\!\left(\frac1n\right)\to0.
\]
因此在 $I$ 上
\[
  S_n^f(x)=S_n^g(x)+S_n^h(x)
\]
一致收敛于 $g(x)=f(x)$。故 $f$ 的 Fourier 级数在任何不含其间断点的闭区间上一致收敛。
\end{xsolution}

\paragraph{参考题 10.}
\begin{xsolution}
由参考题 5(1)，取 $x=1$，立即有
\[
  \boxed{\displaystyle
  \sum_{n=1}^{\infty}\frac{\sin n}{n}=\frac{\pi-1}{2}}.
\]

又由 §15.2 练习题 10 的公式
\[
  \sum_{n=1}^{\infty}\frac{\sin^2na}{n^2}=\frac{a(\pi-a)}2,
  \qquad0<a<\pi,
\]
取 $a=1$，得到
\[
  \boxed{\displaystyle
  \sum_{n=1}^{\infty}\frac{\sin^2n}{n^2}=\frac{\pi-1}{2}}.
\]

最后令
\[
  H(x)=\sum_{n=1}^{\infty}\frac{1-\cos nx}{n^4}.
\]
该级数可逐项求两次导数到
\[
  H''(x)=\sum_{n=1}^{\infty}\frac{\cos nx}{n^2}
  =\frac{3x^2-6\pi x+2\pi^2}{12},
  \qquad0\le x\le2\pi,
\]
并且 $H(0)=H'(0)=0$。所以
\[
  H(x)=\int_0^x(x-t)
  \frac{3t^2-6\pi t+2\pi^2}{12}\,\dd t.
\]
取 $x=2$，直接积分得
\[
  H(2)=\frac{(\pi-1)^2}{3}.
\]
由于 $1-\cos2n=2\sin^2n$，
\[
  2\sum_{n=1}^{\infty}\frac{\sin^2n}{n^4}=H(2),
\]
故
\[
  \boxed{\displaystyle
  \sum_{n=1}^{\infty}\frac{\sin^2n}{n^4}=\frac{(\pi-1)^2}{6}}.
\]
\end{xsolution}

\paragraph{参考题 11.}
\begin{xsolution}
只需证明任意一个长度为 $2\pi$ 的周期区间内至少有 $2n+2$ 次变号。反设某一周期内只有 $2m$ 个变号点
\[
  x_1,x_2,\ldots,x_{2m},\qquad 2m\le2n.
\]
周期函数绕一周后的符号必须回到原符号，所以变号次数为偶数。

考虑
\[
  P(x)=C\prod_{j=1}^{2m}\sin\frac{x-x_j}{2},
\]
其中常数 $C=\pm1$ 选取成使 $f(x)P(x)\ge0$。每跨过一个 $x_j$，$f$ 和 $P$ 都同时改变符号，因此这样的 $C$ 可以统一选取。又因 $f\not\equiv0$ 且连续，必在某个开区间上 $fP>0$，于是
\[
  \int_{-\pi}^{\pi}f(x)P(x)\,\dd x>0.
\]

另一方面，$P$ 是 $2m$ 个半角正弦因子的乘积。把每个因子写成指数形式可见其频率从 $-m$ 到 $m$，故 $P$ 是次数不超过 $m\le n$ 的三角多项式。题设说明 $f$ 与所有次数不超过 $n$ 的三角多项式正交，所以
\[
  \int_{-\pi}^{\pi}f(x)P(x)\,\dd x=0,
\]
矛盾。

故每个完整周期内至少有 $2n+2$ 次变号。任意长度大于 $2\pi$ 的区间包含一个完整周期，因此结论成立。
\end{xsolution}

\paragraph{参考题 12.}
\begin{xsolution}
若 $f$ 单调增加且趋于 $0$，则对 $-f$ 作以下论证即可，所以只考虑 $f\ge0$ 单调减少且 $f(x)\to0$ 的情形。

令
\[
  A_n(x)=\int_0^x\sin nt\,\dd t=\frac{1-\cos nx}{n},
\]
则
\[
  0\le A_n(x)\le\frac2n.
\]
对任意 $R>0$，按 Stieltjes 分部积分（或积分第二中值定理）有
\[
  \int_0^Rf(x)\sin nx\,\dd x
  =f(R)A_n(R)-\int_0^RA_n(x)\,\dd f(x).
\]
因为 $\dd f\le0$，两项均可估计，得到
\[
  \left|\int_0^Rf(x)\sin nx\,\dd x\right|
  \le\frac{2f(0)}n.
\]
令 $R\to+\infty$，由 Dirichlet 判别法广义积分存在，并保持同一估计。因此
\[
  \left|\int_0^{+\infty}f(x)\sin nx\,\dd x\right|
  \le\frac{2\abs{f(0)}}n\to0.
\]
故
\[
  \boxed{\displaystyle
  \lim_{n\to\infty}\int_0^{+\infty}f(x)\sin nx\,\dd x=0 }.
\]
\end{xsolution}

\paragraph{参考题 13.}
\begin{xsolution}
因为
\[
  \sum_{n=1}^{\infty}n^2\ee^{-n}<\infty,
\]
题给三角级数绝对一致收敛，所以其和函数 $f$ 连续，而且由命题 15.1.1，该级数就是 $f$ 的 Fourier 级数。特别地，第一个正弦 Fourier 系数为
\[
  b_1=\ee^{-1}.
\]
设
\[
  M=\max_{0\le x\le2\pi}\abs{f(x)}.
\]
由 Fourier 系数公式
\[
  \frac1\ee
  =\frac1\pi\int_0^{2\pi}f(x)\sin x\,\dd x,
\]
故
\[
  \frac1\ee
  \le\frac{M}{\pi}\int_0^{2\pi}\abs{\sin x}\,\dd x
  =\frac{4M}{\pi}.
\]
于是
\[
  M\ge\frac{\pi}{4\ee}.
\]
由于 $\pi/4>2/\pi$，进一步得到题目所要求的较弱估计
\[
  \boxed{\displaystyle
  M\ge\frac2{\pi\ee}}.
\]
\end{xsolution}

\paragraph{参考题 14.}
\begin{xsolution}
因为 $\varepsilon_n\to0$，可递归选取严格增加的正整数列
\[
  n_1<n_2<\cdots
\]
使
\[
  \varepsilon_{n_k}<2^{-k},\qquad k=1,2,\ldots.
\]
定义
\[
  f(x)=\sum_{k=1}^{\infty}2^{-k}\cos n_kx.
\]
由
\[
  \sum_{k=1}^{\infty}2^{-k}<\infty
\]
和 M-判别法，该级数在 $\RR$ 上绝对一致收敛，所以 $f$ 是周期 $2\pi$ 的连续函数。又由命题 15.1.1，这个一致收敛的三角级数就是 $f$ 的 Fourier 级数。因此
\[
  a_{n_k}=2^{-k},\qquad b_{n_k}=0.
\]
于是对所有 $k$，
\[
  \abs{a_{n_k}}+\abs{b_{n_k}}
  =2^{-k}>\varepsilon_{n_k}.
\]
这样的指标 $n_k$ 有无限多个，结论得证。
\end{xsolution}

\paragraph{参考题 15.}
\begin{xsolution}
记 $f'$ 的 Fourier 系数为 $a_n',b_n'$。首先
\[
  a_0'=\frac1\pi\int_{-\pi}^{\pi}f'(x)\,\dd x
  =\frac{f(\pi)-f(-\pi)}{\pi}=c.
\]
对 $n\ge1$，分部积分得
\begin{align*}
  a_n'
  &=\frac1\pi\int_{-\pi}^{\pi}f'(x)\cos nx\,\dd x\\
  &=\frac{(-1)^n[f(\pi)-f(-\pi)]}{\pi}
    +\frac n\pi\int_{-\pi}^{\pi}f(x)\sin nx\,\dd x\\
  &=(-1)^nc+nb_n,
\end{align*}
以及
\begin{align*}
  b_n'
  &=\frac1\pi\int_{-\pi}^{\pi}f'(x)\sin nx\,\dd x\\
  &=-\frac n\pi\int_{-\pi}^{\pi}f(x)\cos nx\,\dd x
  =-na_n,
\end{align*}
因为边界上的 $\sin n\pi$ 为 $0$。故
\[
  \boxed{\displaystyle
  f'(x)\sim\frac c2+
  \sum_{n=1}^{\infty}
  [(nb_n+(-1)^nc)\cos nx-na_n\sin nx]}.
\]

再由 Riemann 引理，$a_n'\to0$，所以
\[
  nb_n+(-1)^nc\to0.
\]
乘以 $(-1)^{n-1}$，得到
\[
  \boxed{\displaystyle
  c=\lim_{n\to\infty}(-1)^{n-1}nb_n }.
\]
\end{xsolution}

\paragraph{参考题 16.}
\begin{xsolution}
记题设函数 $\varphi$ 的 Fourier 系数为
\[
  A_0=c,\qquad
  A_n=nb_n+(-1)^nc,\qquad
  B_n=-na_n\quad(n\ge1).
\]
对 $\varphi$ 应用 Fourier 级数逐项积分定理。令
\[
  G(x)=\int_0^x\left(\varphi(t)-\frac c2\right)\,\dd t.
\]
则其积分后的 Fourier 级数为
\begin{align*}
  G(x)
  &=\sum_{n=1}^{\infty}\frac{B_n}{n}
   +\sum_{n=1}^{\infty}
   \left(-\frac{B_n}{n}\cos nx+\frac{A_n}{n}\sin nx\right)\\
  &=-\sum_{n=1}^{\infty}a_n
   +\sum_{n=1}^{\infty}
   \left[a_n\cos nx+\left(b_n+\frac{(-1)^nc}{n}\right)\sin nx\right].
\end{align*}
逐项积分定理还保证右端三角级数一致收敛。因此
\[
  H(x):=\sum_{n=1}^{\infty}
  \left[a_n\cos nx+\left(b_n+\frac{(-1)^nc}{n}\right)\sin nx\right]
\]
一致收敛。

另一方面，参考题 6(3) 已给出
\[
  \sum_{n=1}^{\infty}\frac{(-1)^{n-1}}n\sin nx
  =\frac x2,
  \qquad\abs{x}<\pi,
\]
并且该级数在 $x=\pm\pi$ 也收敛（其和为 $0$）。因此原三角级数可写为
\[
  \frac{a_0}{2}+H(x)
  +c\sum_{n=1}^{\infty}\frac{(-1)^{n-1}}n\sin nx,
\]
所以处处收敛。记其和为 $f$。

在 $-\pi<x<\pi$ 内，由上面的 $G$ 表示式可知存在常数 $C$ 使
\[
  f(x)=C+G(x)+\frac c2x.
\]
于是 $f$ 在开区间内绝对连续，并且在 $\varphi$ 的每个连续点，微积分基本定理给出
\[
  f'(x)=G'(x)+\frac c2
  =\varphi(x)-\frac c2+\frac c2
  =\varphi(x).
\]

还需验证给定三角级数确为 $f$ 的 Fourier 级数。由于 $H$ 一致收敛，可以逐项积分求其 Fourier 系数；而
\[
  \sum_{n=1}^{\infty}\frac{(-1)^{n-1}}n\sin nx
\]
本身就是区间 $(-\pi,\pi)$ 上函数 $x/2$ 的 Fourier 级数。因此 $f$ 的非零频率 Fourier 系数正是给定的 $a_n,b_n$。

再看常数项。由上面的分解，原三角级数的和函数就是
\[
  f(x)=\frac{a_0}{2}+H(x)
  +c\sum_{n=1}^{\infty}\frac{(-1)^{n-1}}n\sin nx.
\]
一致收敛级数 $H$ 的常数 Fourier 系数为 $0$，而最后的正弦级数也有零平均值，所以
\[
  \frac1\pi\int_{-\pi}^{\pi}f(x)\,\dd x=a_0.
\]
故 $f$ 的全部 Fourier 系数确为 $a_0,a_n,b_n$，原三角级数就是 $f$ 的 Fourier 级数，并且在 $\varphi$ 的连续点有 $f'=\varphi$。
\end{xsolution}

\paragraph{参考题 17.}
\begin{xsolution}
对题给三角级数令
\[
  a_n=0,
  \qquad b_1=0,
  \qquad b_n=\frac{(-1)^n n}{n^2-1}\quad(n\ge2).
\]
则
\[
  c=\lim_{n\to\infty}(-1)^{n-1}nb_n
  =-1.
\]
按照参考题 16，应考察函数的 Fourier 级数
\[
  \frac c2+
  \sum_{n=1}^{\infty}[nb_n+(-1)^nc]\cos nx.
\]
当 $n=1$ 时系数为 $1$；当 $n\ge2$ 时
\[
  nb_n+(-1)^nc
  =\frac{(-1)^n}{n^2-1}.
\]
故
\[
  \varphi(x)=-\frac12+\cos x+
  \sum_{n=2}^{\infty}\frac{(-1)^n}{n^2-1}\cos nx.
\]
后一个级数绝对一致收敛，所以 $\varphi$ 连续。由参考题 16，原三角级数处处收敛，并且在 $(-\pi,\pi)$ 内其和函数可导，导函数为 $\varphi$；同时该三角级数是这个和函数的 Fourier 级数。

现在求与该 Fourier 级数对应的连续可微代表函数。令
\[
  \boxed{\displaystyle
  F(x)=\frac{x\cos x}{2}+\frac{\sin x}{4},
  \qquad -\pi\le x\le\pi }.
\]
则 $F$ 在 $[-\pi,\pi]$ 上连续可微，并且
\[
  F''(x)+F(x)=-\sin x.
\]
下面核对其 Fourier 系数。$F$ 为奇函数，所以所有余弦系数及常数项均为 $0$。对 $n\ge2$，利用
\[
  \sin nx\cos x=\frac12[\sin(n+1)x+\sin(n-1)x]
\]
以及
\[
  \int_0^\pi x\sin mx\,\dd x=-\frac{\pi(-1)^m}{m},
\]
可得
\[
  b_n=\frac2\pi\int_0^\pi F(x)\sin nx\,\dd x
  =\frac{(-1)^n n}{n^2-1}.
\]
当 $n=1$ 时，$x\cos x/2$ 对 $b_1$ 的贡献为 $-1/4$，而 $\sin x/4$ 的贡献为 $1/4$，故 $b_1=0$。因此 $F$ 的 Fourier 级数正是题给级数。

需要注意，题给三角级数本身在 $x=\pm\pi$ 处每项都为 $0$，故其逐点和在这两个点为 $0$；而连续函数 $F$ 满足
\[
  F(\pi)=-\frac\pi2,\qquad F(-\pi)=\frac\pi2.
\]
单点取值不影响 Fourier 系数，所以题给级数仍是这个连续可微函数 $F$ 的 Fourier 级数；这里要求的是“某个连续可微函数的 Fourier 级数”，并不要求该级数在端点逐点收敛到该函数的端点值。
\end{xsolution}

\paragraph{参考题 18.}
\begin{xsolution}
设给定有限正交系为
\[
  e_1,e_2,\ldots,e_m.
\]
在 $(a,b)$ 内取 $m+1$ 个两两不交的小区间，并在第 $j$ 个小区间内取一个非零连续函数 $\phi_j$，使各 $\phi_j$ 的支集两两不交。

考虑齐次线性方程组
\[
  \sum_{j=1}^{m+1}c_j\int_a^b\phi_j(x)e_i(x)\,\dd x=0,
  \qquad i=1,\ldots,m.
\]
这是 $m$ 个方程、$m+1$ 个未知数，因此存在非零解
$(c_1,\ldots,c_{m+1})$。令
\[
  g(x)=\sum_{j=1}^{m+1}c_j\phi_j(x).
\]
由于各支集互不相交且系数不全为零，$g\not\equiv0$；同时 $g$ 连续，从而可积且平方可积。由方程组定义，
\[
  \int_a^bg(x)e_i(x)\,\dd x=0,
  \qquad i=1,\ldots,m.
\]
所以存在非零平方可积函数与该有限正交系的每个元素正交，故有限正交系不可能完备。
\end{xsolution}

\paragraph{参考题 19（de la Vallée Poussin 核）.}
\begin{xsolution}
记
\[
  K_n(u)=\frac{(2n)!!}{2\pi(2n-1)!!}
  \left(\cos\frac u2\right)^{2n}.
\]
则
\[
  V_n(x)=\int_{-\pi}^{\pi}f(t)K_n(t-x)\,\dd t.
\]

(1) 利用二项式展开可得
\[
  \left(\cos\frac u2\right)^{2n}
  =2^{-2n}\binom{2n}{n}
  +2^{1-2n}\sum_{k=1}^{n}\binom{2n}{n-k}\cos ku.
\]
把 $u=t-x$ 代入，并用
\[
  \cos k(t-x)=\cos kt\cos kx+\sin kt\sin kx,
\]
对 $t$ 积分后，$V_n(x)$ 成为
\[
  A_0+\sum_{k=1}^{n}(A_k\cos kx+B_k\sin kx),
\]
故 $V_n$ 是 $n$ 次三角多项式。

(2) 先验证归一化。由 Wallis 积分公式
\[
  \int_0^{\pi/2}\cos^{2n}v\,\dd v
  =\frac\pi2\frac{(2n-1)!!}{(2n)!!},
\]
所以
\[
  \int_{-\pi}^{\pi}K_n(u)\,\dd u=1.
\]
而且 $K_n\ge0$。

再证明核的质量集中于 $u=0$。给定 $0<\delta<\pi$，在
$\delta\le\abs{u}\le\pi$ 上有
\[
  \abs{\cos(u/2)}\le q:=\cos(\delta/2)<1.
\]
还需控制归一化常数。对 $\abs{u}\le n^{-1/2}$，
\[
  \cos\frac u2\ge1-\frac{u^2}{8}\ge1-\frac1{8n},
\]
故存在与 $n$ 无关的常数 $c_0>0$，使
\[
  \int_{-\pi}^{\pi}\left(\cos\frac u2\right)^{2n}\,\dd u
  \ge\frac{c_0}{\sqrt n}.
\]
由于 $K_n$ 的归一化常数正是上述积分的倒数，故它为 $\bigO(\sqrt n)$。于是
\[
  \int_{\delta\le\abs{u}\le\pi}K_n(u)\,\dd u
  \le C\sqrt n\,q^{2n}\to0.
\]

最后利用 $f$ 的一致连续性。给定 $\varepsilon>0$，先取 $\delta$ 使
\[
  \abs{f(x+u)-f(x)}<\frac\varepsilon2
  \quad(\abs{u}<\delta)
\]
对所有 $x$ 一致成立。设 $M=\norm{f}_\infty$。由周期性和平移变量，
\begin{align*}
  \abs{V_n(x)-f(x)}
  &\le\int_{-\pi}^{\pi}K_n(u)
  \abs{f(x+u)-f(x)}\,\dd u\\
  &\le\frac\varepsilon2
  +2M\int_{\delta\le\abs{u}\le\pi}K_n(u)\,\dd u.
\end{align*}
第二项当 $n$ 足够大时小于 $\varepsilon/2$，并且估计与 $x$ 无关。因此
\[
  \sup_{x\in\RR}\abs{V_n(x)-f(x)}\to0,
\]
即 $V_n$ 在 $(-\infty,+\infty)$ 上一致收敛于 $f$。
\end{xsolution}

\paragraph{参考题 20.}
\begin{xsolution}
记
\[
  S_n(x)=\sum_{k=1}^{n}\frac{\cos kx}{k}.
\]
先证明 $S_n(x)\ge-1$。

我们使用一个简单的 Fejér 正性引理。设
$d_0,d_1,\ldots,d_N$ 满足
$d_{N+1}=d_{N+2}=0$ 且
\[
  \Delta^2d_m:=d_m-2d_{m+1}+d_{m+2}\ge0.
\]
标准 Fejér 核
\[
  K_m(x)=1+2\sum_{k=1}^{m}\left(1-\frac{k}{m+1}\right)\cos kx
  =\frac1{m+1}\left(\frac{\sin\frac{m+1}{2}x}{\sin\frac x2}\right)^2
\]
非负。比较 Fourier 系数可验证恒等式
\[
  d_0+2\sum_{k=1}^{N}d_k\cos kx
  =\sum_{m=0}^{N}(m+1)\Delta^2d_m\,K_m(x),
\]
故左边非负。

现在
\[
  S_n'(x)=-\sum_{k=1}^{n}\sin kx
  =-\frac{\sin(nx/2)\sin((n+1)x/2)}{\sin(x/2)}.
\]
因此在一个周期内，除端点外的极值点只有两类：
\[
  x=\frac{2\pi j}{n}
  \quad\text{或}\quad
  x=\frac{2\pi j}{n+1}.
\]

若 $x=2\pi j/n$ 为非端点，则
\[
  \sum_{k=1}^{n-1}\cos kx=-1.
\]
令
\[
  d_0=1,\qquad
  d_k=\frac12\left(\frac1k-\frac1n\right)\quad(1\le k\le n-1).
\]
数列 $1/k$ 为凸数列，且端点处的二阶差分也非负，所以由上面的 Fejér 正性引理，
\[
  1+\sum_{k=1}^{n-1}\left(\frac1k-\frac1n\right)\cos kx\ge0.
\]
利用根单位和式，左边恰为 $1+S_n(x)$，故 $S_n(x)\ge-1$。

若 $x=2\pi j/(n+1)$，则
\[
  \sum_{k=1}^{n}\cos kx=-1.
\]
这次令
\[
  d_0=\frac{n}{n+1},\qquad
  d_k=\frac12\left(\frac1k-\frac1{n+1}\right)\quad(1\le k\le n).
\]
其二阶差分同样全部非负，因而
\[
  \frac{n}{n+1}+
  \sum_{k=1}^{n}\left(\frac1k-\frac1{n+1}\right)\cos kx\ge0.
\]
利用 $\sum_{k=1}^n\cos kx=-1$，左边仍恰为 $1+S_n(x)$。所以所有内部极值点都有 $S_n\ge-1$；端点 $x=0,2\pi$ 处 $S_n=\sum_{k=1}^n1/k>0$。故
\[
  \boxed{S_n(x)\ge-1\quad\text{对一切 }x}.
\]

下面求最小值的极限。记
\[
  m_n=\min_{0\le x\le\pi}S_n(x).
\]
在 $x=\pi$，
\[
  S_n(\pi)=\sum_{k=1}^{n}\frac{(-1)^k}{k}\longrightarrow-\ln2,
\]
所以
\[
  \limsup_{n\to\infty}m_n\le-\ln2.
\]

另一方面，对 $0<x<2\pi$，由 Abel 极限定理
\[
  \sum_{k=1}^{\infty}\frac{\cos kx}{k}
  =-\ln\left(2\sin\frac x2\right)=:L(x).
\]
在每个 $[\delta,\pi]$ 上，由 Dirichlet 判别法余项估计，$S_n\to L$ 一致。

只需排除最小点趋向 $0$。若 $x_n\to0$ 且 $nx_n$ 有界，则由
$1-\cos y\le y^2/2$，
\begin{align*}
  S_n(x_n)
  &=\sum_{k=1}^{n}\frac1k
    -\sum_{k=1}^{n}\frac{1-\cos(kx_n)}k\\
  &\ge H_n-\frac{x_n^2}{2}\sum_{k=1}^{n}k
  =H_n-\bigO(1)\to+\infty.
\end{align*}
若 $x_n\to0$ 而 $nx_n\to\infty$，则 Dirichlet 余项估计给出
\[
  \left|L(x_n)-S_n(x_n)\right|
  \le\frac{C}{(n+1)\sin(x_n/2)}
  \le\frac{C'}{nx_n}\to0,
\]
而 $L(x_n)\to+\infty$，故仍有 $S_n(x_n)\to+\infty$。因此取得最小值的点不可能趋向 $0$。

任取实现 $\liminf m_n$ 的子列及其最小点，再取收敛子列 $x_n\to x_*>0$。在包含 $x_*$ 的闭区间上有一致收敛，故
\[
  \liminf_{n\to\infty}m_n=L(x_*)
  \ge\min_{0<x\le\pi}L(x)=-\ln2.
\]
结合上面的 $\limsup$，得到
\[
  \boxed{\displaystyle
  \lim_{n\to\infty}\min_{0\le x\le\pi}S_n(x)=-\ln2 }.
\]
\end{xsolution}
